\chapter{总结与展望}
\label{chap:conclusion}

\section{工作总结}

随着 B5G/6G 网络向海量连接与多样化业务需求持续演进，传统随机接入机制在复杂物联网场景中的性能瓶颈日益凸显。针对重复竞争开销大、系统效率与公平性难以兼顾、动态场景适配能力不足等问题，全文按照单接入机制理论优化到智能接入协议决策的逻辑主线，展开系统研究。主要工作可归纳如下：
\begin{itemize}

\item 在授权频谱随机接入场景中，针对短报文按包竞争引起的碰撞累积与信令开销问题，提出连续传输增强 Aloha 机制 SAST。理论分析与仿真结果表明，该机制在不改变原有分布式接入结构的前提下，将系统最大数据吞吐量提升至0.5，突破了经典时隙 Aloha 的性能上限$1/e$。
进一步地，所建立的分析结果揭示了连续传输策略通过延长信道占用周期，有效分摊竞争开销的内在机理，使系统吞吐量与接入效率得到同步提升。同时，在5G小数据传输场景中的对比分析表明，相较于两步随机接入方案，该机制在提升吞吐量的同时显著降低信令开销，在饱和条件下可降低约一半，在非饱和条件下降幅可达70\%。
上述结果表明，连续传输机制在维持现有协议架构兼容性的前提下，能够有效提升海量物联网场景下的接入效率与资源利用率，为5G及后续系统中高效随机接入机制的设计提供了理论依据与实现路径。


\item 在非授权频谱场景中，针对 LBT 机制下频繁侦听与退避导致的信道利用率下降问题，本文将连续传输思想引入载波侦听多路访问机制，提出连续传输增强 CSMA 协议 CSMAST。
该机制在保留载波侦听竞争基本流程的基础上，通过在成功接入后支持连续数据传输，有效降低重复竞争带来的开销。考虑空闲侦听、碰撞与成功传输对应时隙长度异构的特性，本文构建了基于吸收马尔可夫链的休假排队模型，对信道假期过程进行刻画，并推导出系统吞吐量闭环数学表达式。在此基础上，引入并且成功推导 Jain 公平性指标 $J_T$的表达式，分析了连续传输机制对系统吞吐量和公平性的影响。理论与仿真结果表明，CSMAST 在同等侦听条件下，可显著提升系统吞吐量大约28.8\%，同时能够在提升信道利用率的同时保持良好的公平性。此外，结果进一步揭示了竞争接入系统中吞吐量与公平性之间的内在权衡关系。

\textcolor{red}{
    \item 在动态业务环境场景中，针对单一接入机制及参数配置在动态业务变化中适配能力不足的问题，构建统一的接入评估与决策框架。该框架通过离散事件仿真对不同接入协议及参数配置在多种业务场景下的性能进行系统采样，形成覆盖多维性能指标的高质量数据集。在此基础上，本文进一步构建多任务学习模型，对吞吐量、时延、公平性及信令开销等关键性能进行联合预测。结果表明，该模型能够准确刻画复杂接入环境中的性能映射关系，实现对多协议性能的高精度快速估计。进一步地，基于所构建的预测模型与效用函数 $U(\tilde{\mathbf{y}})$，设计面向动态业务环境的自适应决策机制，实现协议选择与参数配置的联合优化。实验结果表明，所提出方法能够在保证良好实时性的情况下，获得业务动态变化下获得稳定的性能增益。}
\end{itemize}


\section{研究展望}
在上述研究中，本文围绕随机接入机制的性能优化与动态环境自适应决策问题，分别从单机制增强与多机制协同决策两个层面开展了系统研究，在提升接入效率、降低竞争开销以及实现多性能指标协同优化方面取得了一定进展。然而，随着 B5G/6G 网络向更高密度连接、更复杂业务形态及更动态网络环境持续演进，接入系统在物理信道特性、业务异构性以及系统规模等方面仍面临更加复杂的挑战，现有研究在建模精细性、决策机制鲁棒性及工程实现等方面仍存在进一步拓展空间。基于此，未来可从以下几个方向展开深入研究：
\begin{itemize}
\item 在物理层建模方面，现有分析基于理想碰撞信道模型，假设冲突数据包完全不可恢复，难以刻画实际无线信道中存在的捕获效应与干扰结构。未来可引入多包接收（Multi-Packet Reception, MPR）与干扰消除（Interference Cancellation, IC）等机制，对部分可解码冲突场景进行建模，从而建立跨层耦合接入模型。在此基础上，有望重新刻画本文所提的增强机制的性能上界与稳定性区域，进一步挖掘连续传输机制在复杂信道环境下的性能潜力。

\item 在网络模型方面，现有研究主要基于同质节点假设，未充分考虑多业务类型在时延、可靠性及吞吐需求上的差异。在实际物联网系统中，往往存在不同优先级的异构设备共存的情况。未来可面向异构设备接入场景，构建区分业务类型的随机接入模型，引入优先级机制与差异化接入策略，研究不同业务在共享信道中的资源竞争与性能保障问题，并探索面向服务质量（QoS）约束的异构节点接入机制设计方法。

\item 在智能决策方面，现有方法基于监督学习对接入系统的性能映射关系进行建模，并在此基础上实现协议与参数的自适应决策。然而，该方法依赖于离线仿真数据进行训练，在面对分布发生变化的超复杂动态环境时，其泛化能力与适应能力仍有待提升。未来可进一步研究面向动态环境的在线学习与增量学习方法，使模型能够在运行过程中持续更新，从而提升对未知业务场景的适应能力。此外，现有模型依赖集中式训练与全局数据，在大规模网络中可能带来较高的数据获取与计算开销。未来可探索轻量化模型设计与分布式学习机制，在保证预测精度的同时降低计算复杂度，并使节点能够基于局部信息进行近似决策，从而提升系统的可扩展性与实际部署可行性。

\item 在工程实现方面，现有研究主要基于理论分析与离散事件仿真验证，尚未考虑实际商用系统中的实现约束与真实协议开销。本文所提出的连续传输增强机制在设计上保持了对现有 3GPP 随机接入协议框架的高度兼容性，仅在接入成功后的传输阶段进行扩展，因此具备较低的协议改动成本与良好的工程落地潜力。未来可在此基础上，结合实际系统约束，进一步研究机制在标准协议栈中的实现方式与信令交互过程，并基于软件无线电（Software Defined Radio, SDR）平台开展原型系统设计与实验验证，评估其在真实无线环境中的性能表现、实现复杂度及系统开销，从而推动相关机制向实际通信系统中的应用转化。
\end{itemize}





